\documentclass[12pt,a4paper]{article}

\usepackage[utf8]{inputenc}
\usepackage[T1]{fontenc}
\usepackage[english]{babel}
\usepackage{geometry}
\usepackage{hyperref}
\usepackage{listings}
\usepackage{xcolor}
\usepackage{booktabs}
\usepackage{parskip}
\usepackage{enumitem}
\usepackage{titlesec}
\usepackage{fancyhdr}
\usepackage{amsmath}
\usepackage{textcomp}   % FIX: needed for \textperiodcentered and safe symbol mapping

\geometry{margin=2.5cm}

\hypersetup{colorlinks=true, linkcolor=blue, urlcolor=blue}

\definecolor{codebg}{HTML}{1a1d27}
\definecolor{codefg}{HTML}{e2e8f0}
\definecolor{codecomment}{HTML}{94a3b8}
\definecolor{codestring}{HTML}{f59e0b}
\definecolor{codekw}{HTML}{4f7ef8}

\lstset{
    basicstyle=\ttfamily\small\color{codefg},
    backgroundcolor=\color{codebg},
    frame=single,
    rulecolor=\color{codebg},
    breaklines=true,
    keywordstyle=\color{codekw}\bfseries,
    commentstyle=\color{codecomment}\itshape,
    stringstyle=\color{codestring},
    language=Python,
    showstringspaces=false,
    xleftmargin=6pt, xrightmargin=6pt,
    aboveskip=8pt, belowskip=8pt
}

\pagestyle{fancy}
\fancyhf{}
\rhead{TP02 v2 --- Secure Messaging}   % FIX: --- instead of Unicode em dash —
\lhead{Denis Ozdemir --- ICE3 ENSIBS} % FIX: idem
\rfoot{\thepage}

\titleformat{\section}{\large\bfseries}{{\S}\,\thesection}{1em}{}
\titleformat{\subsection}{\normalsize\bfseries}{\thesubsection}{1em}{}

% ──────────────────────────────────────────────────────────────────────────────
\title{%
  \vspace{-1cm}
  \textbf{Advanced Secure Messaging}\\[6pt]
  \large AES-GCM $\cdot$ RSA-PSS Signatures $\cdot$ Multi-session $\cdot$ Audit Log
  % FIX: $\cdot$ instead of Unicode middle dot ·  (U+00B7, not in base T1 mapping)
}
\author{Denis O.}
\date{\today}

\begin{document}

\maketitle
\tableofcontents

% ──────────────────────────────────────────────────────────────────────────────
\section{Objective and Scope}

Version~2 extends the baseline v1 system with production-grade security properties:

\begin{itemize}[noitemsep]
  \item \textbf{Authenticated encryption}: AES-CBC (unauthenticated) is replaced by AES-GCM, which simultaneously encrypts and provides a 128-bit integrity tag --- preventing ciphertext tampering without a separate HMAC.
  \item \textbf{Non-repudiation}: every outgoing message is signed with the sender's RSA private key (PSS/SHA-256). Recipients verify the signature against the sender's registered public key.
  \item \textbf{Multi-session}: each client maintains a per-peer AES key dictionary instead of a single global variable, enabling simultaneous conversations.
  \item \textbf{Server-side signature verification}: the server verifies the session-request signature before issuing a new AES key, preventing impersonation.
  \item \textbf{Idempotent session management}: duplicate session requests return HTTP~409 rather than silently re-generating a key; a \texttt{/session/reset} endpoint allows intentional re-keying.
  \item \textbf{Structured messages}: message payloads carry \texttt{text}, \texttt{timestamp}, \texttt{signature}, \texttt{sender}, and \texttt{session\_id} --- enabling richer UI and auditability.
  \item \textbf{Audit log}: the server persists a timestamped log of all cryptographic events accessible via \texttt{GET /audit}.
  \item \textbf{Modern UI}: a dark-mode interface with message timestamps, RSA-PSS verification badges, encryption summary panel, and toast notifications.
\end{itemize}

% ──────────────────────────────────────────────────────────────────────────────
\section{Architecture}

\subsection{Component overview}

\begin{center}
\begin{tabular}{llp{8cm}}
\toprule
\textbf{Process} & \textbf{Port} & \textbf{Role} \\
\midrule
\texttt{server.py} & 5000 & Key distribution, signature verification, audit log \\
\texttt{client.py} Alice & 5001 & Chat client, AES-GCM encrypt/decrypt, RSA-PSS sign/verify \\
\texttt{client.py} Bob & 5002 & Chat client (same code, different port/username) \\
\bottomrule
\end{tabular}
\end{center}

\subsection{Full message flow}

\begin{enumerate}[noitemsep]
  \item Both clients generate RSA-2048 key pairs locally at startup.
  \item Each client \texttt{POST /register} their public key with the server.
  \item Alice calls \texttt{POST /session} with a PSS-signed proof-of-identity string \texttt{"Alice:Bob"}.
  \item Server verifies the signature, generates a 16-byte AES key, encrypts it under each party's RSA public key, records the session in the audit log, and returns both ciphertexts.
  \item Alice decrypts her copy; forwards Bob's RSA-wrapped copy (plus session ID) via peer-to-peer \texttt{POST /receive\_key}.
  \item Alice composes message payload \(\{\texttt{text},\texttt{timestamp},\texttt{signature},\texttt{session\_id}\}\), signs the plaintext with RSA-PSS, serialises to JSON, and encrypts the whole JSON string with AES-GCM.
  \item The AES-GCM blob is sent to Bob via \texttt{POST /receive\_message}.
  \item Bob decrypts with AES-GCM (tag verification happens automatically), parses the JSON, fetches Alice's public key, and verifies the RSA-PSS signature.
  \item Bob's UI displays the message with a green badge if verification succeeds.
\end{enumerate}

% ──────────────────────────────────────────────────────────────────────────────
\section{Cryptographic Details}

\subsection{AES-GCM authenticated encryption}

AES-128-GCM provides \emph{confidentiality} (via CTR-mode encryption) and \emph{integrity} (via GHASH authentication tag) in a single pass.  The 96-bit nonce is generated randomly for each message and prepended to the tag+ciphertext blob before base64 encoding.

\[
  \text{blob} = \underbrace{\text{nonce}}_{16\,\text{B}} \;\|\;
                \underbrace{\text{tag}}_{16\,\text{B}} \;\|\;
                \underbrace{\text{ciphertext}}_{|\text{plaintext}|\,\text{B}}
\]

\begin{lstlisting}
def encrypt_gcm(plaintext: str, key: bytes) -> str:
    cipher = AES.new(key, AES.MODE_GCM)
    ciphertext, tag = cipher.encrypt_and_digest(plaintext.encode())
    blob = cipher.nonce + tag + ciphertext
    return base64.b64encode(blob).decode()
\end{lstlisting}

If the tag does not match on decryption, \texttt{decrypt\_and\_verify} raises \texttt{ValueError} --- the message is rejected before any plaintext is exposed.

\subsection{RSA-PSS message signatures}

Each message is signed over its plaintext with the sender's RSA-2048 private key using the PSS scheme and SHA-256:

\begin{lstlisting}
def sign_message(text: str) -> str:
    h   = SHA256.new(text.encode())
    sig = pss.new(my_key).sign(h)
    return base64.b64encode(sig).decode()
\end{lstlisting}

The signature is embedded in the JSON payload \emph{before} encryption, so only the intended recipient (who decrypts successfully) can read and verify it.

\subsection{Server-side session-request signature}

Before issuing an AES session key, the server verifies that the claimed initiator signed the string \(\texttt{"initiator:target"}\):

\begin{lstlisting}
verifier = pss.new(RSA.import_key(clients_db[initiator]["public_key"]))
h        = SHA256.new(f"{initiator}:{target}".encode())
verifier.verify(h, base64.b64decode(sig_b64))   # raises on failure
\end{lstlisting}

This ensures the server does not issue session keys to an entity that cannot prove ownership of the registered private key.

\subsection{Comparison: v1 vs v2}

\begin{center}
\begin{tabular}{lll}
\toprule
\textbf{Property} & \textbf{v1} & \textbf{v2} \\
\midrule
Symmetric cipher     & AES-128-CBC          & AES-128-GCM \\
Integrity            & None (unauthenticated) & 128-bit GCM tag \\
Message signatures   & None                 & RSA-2048 PSS/SHA-256 \\
Session-req.\ authN  & None                 & RSA-PSS proof-of-identity \\
Multi-session        & No (global variable) & Yes (per-peer dict) \\
Idempotent sessions  & No                   & Yes (HTTP 409 + reset) \\
Structured messages  & Plain strings        & JSON with timestamp + sig \\
Audit log            & None                 & Timestamped server-side log \\
HTTP error codes     & 200 for errors       & 400 / 403 / 404 / 409 \\
UI                   & Basic light theme    & Dark theme + sig badges + status \\
\bottomrule
\end{tabular}
\end{center}

% ──────────────────────────────────────────────────────────────────────────────
\section{API Reference}

\subsection{Server endpoints}

\begin{center}
\begin{tabular}{llp{8cm}}
\toprule
Method & Route & Description \\
\midrule
GET  & \texttt{/}                         & Health check + server stats \\
POST & \texttt{/register}                 & Register username + RSA public key \\
GET  & \texttt{/users}                    & List all registered clients \\
POST & \texttt{/session}                  & Issue AES session key (with signature verification) \\
POST & \texttt{/session/reset}            & Force-delete a session between two users \\
GET  & \texttt{/session\_info/<id>}       & Return session metadata \\
GET  & \texttt{/audit}                    & Return full audit event log \\
\bottomrule
\end{tabular}
\end{center}

\subsection{Client endpoints}

\begin{center}
\begin{tabular}{llp{7cm}}
\toprule
Method & Route & Description \\
\midrule
GET  & \texttt{/}                  & Serve the v2 chat UI \\
POST & \texttt{/send}              & Sign, encrypt (GCM), forward message \\
POST & \texttt{/receive\_key}      & Receive RSA-wrapped AES key \\
POST & \texttt{/receive\_key\_v2}  & Extended: also caches sender's public key \\
POST & \texttt{/receive\_message}  & Decrypt (GCM), verify signature, store \\
GET  & \texttt{/messages}          & Return history (\texttt{?peer=X} to filter) \\
GET  & \texttt{/status}            & Expose active sessions + message counts \\
POST & \texttt{/reset\_session}    & Clear local AES key for a peer \\
\bottomrule
\end{tabular}
\end{center}

% ──────────────────────────────────────────────────────────────────────────────
\section{How to Run}

\begin{lstlisting}[language=bash, basicstyle=\ttfamily\small\color{codefg}]
# Terminal 1 - key server
python server.py

# Terminal 2 - Alice
python client.py 5001 Alice

# Terminal 3 - Bob
python client.py 5002 Bob
\end{lstlisting}

Open \url{http://127.0.0.1:5001} and \url{http://127.0.0.1:5002} in two browser tabs. The UI auto-configures the target peer.

To inspect the server audit log:
\begin{lstlisting}[language=bash, basicstyle=\ttfamily\small\color{codefg}]
curl http://127.0.0.1:5000/audit
\end{lstlisting}
% FIX: removed Markdown link syntax [url](url) — inside lstlisting it breaks intent

To view client session status:
\begin{lstlisting}[language=bash, basicstyle=\ttfamily\small\color{codefg}]
curl http://127.0.0.1:5001/status
\end{lstlisting}
% FIX: idem

% ──────────────────────────────────────────────────────────────────────────────
\section{Design Notes and Remaining Limitations}

\begin{itemize}[noitemsep]
  \item \textbf{Server trust}: The server still sees the AES key in plaintext at generation time. A Diffie--Hellman or ECDH key agreement would eliminate this.
  \item \textbf{No TLS}: HTTP transport is unencrypted. Production deployment requires HTTPS.
  \item \textbf{Public key retrieval}: Signature verification relies on fetching the sender's public key from the server via \texttt{GET /users}. In a real system this should be done over a separate authenticated channel, or keys should be certified by a CA.
  \item \textbf{Replay attacks}: messages carry a timestamp but there is no nonce registry to detect replays. Adding a per-session message counter would address this.
  \item \textbf{In-memory storage}: all state is lost on restart. A real deployment would persist sessions and message history in a database.
\end{itemize}

% ──────────────────────────────────────────────────────────────────────────────
\section{Possible Extensions}

\begin{itemize}[noitemsep]
  \item Replace server AES distribution with an ECDH key agreement (zero server trust).
  \item Add TLS to all Flask endpoints (\texttt{ssl\_context} in \texttt{app.run}).
  \item Implement a replay-prevention nonce registry.
  \item Persist sessions and history to SQLite.
  \item Group chat support (one AES key per group, distributed to \(n\) members).
\end{itemize}

% ──────────────────────────────────────────────────────────────────────────────
\section{References}

\begin{itemize}[noitemsep]
  \item NIST SP 800-38D --- Galois/Counter Mode (GCM): \url{https://nvlpubs.nist.gov/nistpubs/Legacy/SP/nistspecialpublication800-38d.pdf}
  \item PKCS \#1 v2.2 / RSA-PSS --- RFC 8017: \url{https://datatracker.ietf.org/doc/html/rfc8017}
  \item PyCryptodome documentation: \url{https://pycryptodome.readthedocs.io}
  \item Flask documentation: \url{https://flask.palletsprojects.com}
\end{itemize}

\end{document}
